\section{Method}
\label{sec:method}



\begin{figure*}[t]
    \centering
    \includegraphics[width=\textwidth]{assets/method.png}
    \caption{Overview of BiCAT. Top: mirror canonicalization $\mathcal{C}$
    flips the signed joints of the right arm ($q_1,\ldots,q_7$, with $q_7$ the
    gripper) and applies shared normalization; the shared encoder $E_{\phi}$
    and causal cross-attention give each arm its partner's context before FSQ;
    and the joint decoder $D_{\theta}$ reconstructs the canonical chunk under
    $\mathcal{L}_{\mathrm{rec}}$, which $\mathcal{C}^{-1}$ maps back to
    $\widehat{\mathbf{A}}$. Bottom: a VLM predicts $\mathbf{y}^{L}$, then
    $\mathbf{y}^{R}$; the frozen decoder returns actions, executed on
    Bimanual SO-101, Unitree G1, and Bimanual YAM. Teal: left arm; orange:
    right arm.}
    \label{fig:bicat_overview}
\end{figure*}


\textbf{BiCAT} tokenizes a bimanual action chunk into two arm-indexed token
streams that share one motion representation and are computed with access to
each other (Fig.~\ref{fig:bicat_overview}). Sec.~\ref{sec:preliminaries}
introduces action tokenization in the bimanual setting,
Sec.~\ref{sec:bicat_tokenization} describes the tokenizer, whose training is
summarized in Algorithm~\ref{alg:bicat_training}, and
Sec.~\ref{sec:bicat_policy} describes how a VLA predicts its tokens.

\subsection{Preliminaries}
\label{sec:preliminaries}

\paragraph{Action tokenization.}
A VLA policy maps an observation $\mathbf{o}$, comprising camera images,
robot state, and a language instruction, to a chunk of $H$ future actions
$\mathbf{A}\in\mathbb{R}^{H\times D}$. To retain the next-token prediction
of its backbone, an action tokenizer $\mathcal{T}$ maps the chunk to discrete
tokens $\mathbf{y}=\mathcal{T}(\mathbf{A})$ from a vocabulary of size $K$,
and a detokenizer $\mathcal{T}^{-1}$ maps tokens back to an executable chunk
$\widehat{\mathbf{A}}=\mathcal{T}^{-1}(\mathbf{y})$. Grouping the tokens into
$B$ blocks $\mathbf{y}_{1},\ldots,\mathbf{y}_{B}$, the policy with
parameters $\theta_{\mathrm{V}}$ is trained on demonstrations $\mathcal{D}$
to maximize
\begin{equation}
    \mathbb{E}_{(\mathbf{o},\mathbf{A})\sim\mathcal{D}}
    \Big[\sum_{b=1}^{B}\log p_{\theta_{\mathrm{V}}}
    (\mathbf{y}_{b}\mid\mathbf{o},\mathbf{y}_{<b})\Big],
    \label{eq:ar_objective}
\end{equation}
where the tokens within a block are predicted in parallel. Token-wise
autoregression uses blocks of one token~\cite{kim2024openvla,pertsch2025fast},
and block-wise prediction~\cite{stern2018blockwise,liu2026oat_rss,jawaid2026nac,dong2026actioncodec,liu2026faster}
reduces the number of sequential decoding steps to $B$.

\paragraph{Bimanual action chunks.}
A bimanual chunk $\mathbf{A}=[\mathbf{A}^{L}\Vert\mathbf{A}^{R}]
\in\mathbb{R}^{H\times 2d_a}$ concatenates the actions of the left and right
arms, where $d_a$ is the per-arm action dimension and $\Vert$ denotes feature
concatenation. Joint tokenization applies one tokenizer to the whole chunk,
$\mathbf{y}=\mathcal{T}(\mathbf{A})$, so that a token can mix both arms.
Per-arm tokenization applies it to each arm separately,
$\mathbf{y}^{r}=\mathcal{T}(\mathbf{A}^{r})$ for $r\in\{L,R\}$, which yields
arm-indexed tokens but computes each arm's tokens without its partner. \textbf{BiCAT}
keeps arm-indexed token streams but computes them from both arms:
\begin{equation}
    (\mathbf{y}^{L},\mathbf{y}^{R})=\mathcal{T}(\mathbf{A}^{L},\mathbf{A}^{R}),
    \qquad
    \widehat{\mathbf{A}}=\mathcal{T}^{-1}
    (\mathbf{y}^{L},\mathbf{y}^{R}),
    \label{eq:bicat_mapping}
\end{equation}
where $\mathbf{y}^{r}\in\{0,\ldots,K-1\}^{M}$ contains $M$ tokens for arm
$r$, and the policy predicts the two streams as $B=2$ successive blocks
(Sec.~\ref{sec:bicat_policy}).

\subsection{Bimanual Coordination Action Tokenization}
\label{sec:bicat_tokenization}




% \paragraph{Canonical action space.}
% Sharing an encoder requires corresponding action channels to use a consistent coordinate convention. We retain the left-arm coordinates and transform the right arm with an embodiment-specific sign vector \(\mathbf{s}\):
% \[
% \overline{\mathbf{A}}^{L}=\mathbf{A}^{L},
% \qquad
% \overline{\mathbf{A}}^{R}=\mathbf{A}^{R}\odot\mathbf{s}.
% \]
% For the YAM platform, we use \(\mathbf{s}=(-1,1,1,1,-1,-1,1)\), where each arm has six joint coordinates and one gripper coordinate. Both canonical arms share the same channel-wise normalization bounds \(\boldsymbol{\ell},\mathbf{u}\in\mathbb{R}^{d_a}\), computed by pooling corresponding coordinates across arms:
% \[
% \mathbf{X}^{r}
% =2\frac{\overline{\mathbf{A}}^{r}-\boldsymbol{\ell}}
% {\mathbf{u}-\boldsymbol{\ell}}-1,
% \qquad r\in\{L,R\}.
% \]
% All operations are element-wise, and the normalization bounds are stored with the tokenizer. This common coordinate system and scaling allow both arms to share encoder parameters.

% \paragraph{Canonical action space.}
% To share motion representations across arms, we first align corresponding
% action channels in a common coordinate convention. We retain the left-arm
% coordinates and transform the right arm with an embodiment-specific sign
% vector $\mathbf{s}\in\{-1,+1\}^{d_a}$:
% \begin{equation}
%     \overline{\mathbf{A}}^{L}=\mathbf{A}^{L},
%     \qquad
%     \overline{\mathbf{A}}^{R}=\mathbf{A}^{R}\odot\mathbf{s},
%     \label{eq:bicat_canonicalization}
% \end{equation}
% where the sign vector is broadcast over time. For YAM, each arm has six joint
% coordinates and one gripper coordinate, and we use
% $\mathbf{s}=(-1,1,1,1,-1,-1,1)$. We then normalize both canonical arms with
% shared channel-wise bounds $\boldsymbol{\ell},\mathbf{u}
% \in\mathbb{R}^{d_a}$, computed by pooling corresponding action channels:
% \begin{equation}
%     \mathbf{X}^{r}
%     =2\frac{\overline{\mathbf{A}}^{r}-\boldsymbol{\ell}}
%     {\mathbf{u}-\boldsymbol{\ell}}-1,
%     \qquad r\in\{L,R\}.
%     \label{eq:bicat_normalization}
% \end{equation}
% All operations are element-wise, and the bounds are stored with the
% tokenizer. The common convention and scaling place corresponding motions in
% the same feature space before they enter the shared encoder.




\paragraph{Canonical action space.}
To share motion representations across arms, we first align corresponding
action channels in a common coordinate convention. We retain the left-arm
coordinates and transform the right arm with an embodiment-specific sign
vector
$\mathbf{s}\in\{-1,+1\}^{d_a}$:
\begin{equation}
    \overline{\mathbf{A}}^{L}=\mathbf{A}^{L},
    \qquad
    \overline{\mathbf{A}}^{R}=\mathbf{A}^{R}\odot\mathbf{s},
    \label{eq:bicat_canonicalization}
\end{equation}
where the sign vector is broadcast over time. Each arm has its joint
coordinates followed by one gripper coordinate, and we use
$\mathbf{s}=(-1,1,1,1,-1,-1,1)$ for the RoboTwin robot and Bimanual YAM,
$\mathbf{s}=(1,-1,-1,1,-1,1,-1,1)$ for Unitree G1, and
$\mathbf{s}=(-1,1,1,1,-1,1)$ for Bimanual SO-101. 
Rather than normalizing each arm independently, we compute shared
channel-wise bounds $\boldsymbol{\ell},\mathbf{u}
\in\mathbb{R}^{d_a}$ by pooling corresponding action channels
across both canonical arms:
\begin{equation}
    \mathbf{X}^{r}
    =2\frac{\overline{\mathbf{A}}^{r}-\boldsymbol{\ell}}
    {\mathbf{u}-\boldsymbol{\ell}}-1,
    \qquad r\in\{L,R\}.
    \label{eq:bicat_normalization}
\end{equation}
All operations are element-wise, and the bounds are stored with the
tokenizer. This shared normalization ensures that corresponding
joint values are mapped identically across arms, aligning their
inputs before the shared encoder. We denote the sign transform followed by
shared normalization as the canonicalization $\mathcal{C}$, and its inverse
as $\mathcal{C}^{-1}$.


\begin{algorithm}[h]
\caption{BiCAT Tokenizer Training}
\label{alg:bicat_training}
\footnotesize
\begin{algorithmic}[1]
\Require Bimanual dataset $\mathcal{D}$;
canonicalization $\mathcal{C}$;
shared encoder $E_{\phi}$;
causal attention modules $\mathrm{SA},\mathrm{CA}$;
projection $P$; quantizer $\mathrm{FSQ}$;
joint decoder $D_{\theta}$.
\While{not converged}
    \State Sample $(\mathbf{A}^{L},\mathbf{A}^{R})
    \sim \mathcal{D}$
    \State \textbf{Canonicalization:}
    $(\mathbf{X}^{L},\mathbf{X}^{R})
    \gets \mathcal{C}(\mathbf{A}^{L},\mathbf{A}^{R})$
    \State \textbf{Encoding:}
    $\mathbf{H}^{L/R}
    \gets E_{\phi}(\mathbf{X}^{L/R})$
    \State \textbf{Coordination modeling:}
    $\mathbf{S}^{L/R}
    \gets \mathrm{SA}(\mathbf{H}^{L/R})$
    \Statex \hspace{\algorithmicindent}
    $\mathbf{C}^{L/R}
    \gets \mathbf{S}^{L/R}
    + \mathrm{CA}(\mathbf{S}^{L/R},
    \mathbf{S}^{R/L})$
    \State \textbf{Quantization:}
    $\hat{\mathbf{Z}}^{L/R}
    \gets \mathrm{FSQ}(P(\mathbf{C}^{L/R}))$
    \State \textbf{Temporal pairing:}
    $\mathbf{M}_i
    \gets \mathrm{Pair}(
    \hat{\mathbf{Z}}_i^{L},
    \hat{\mathbf{Z}}_i^{R})$
    \State \textbf{Joint decoding:}
    $(\hat{\mathbf{X}}^{L},\hat{\mathbf{X}}^{R})
    \gets D_{\theta}(\mathbf{M})$
    \State \textbf{Reconstruction loss:}
    $\mathcal{L}_{\mathrm{rec}}
    \gets \frac{1}{2Hd_a}
    \sum_{r\in\{L,R\}}
    \|\hat{\mathbf{X}}^{r}
    -\mathbf{X}^{r}\|_F^2$
    \State \textbf{Optimization:}
    $\Theta \gets \Theta
    -\eta\nabla_{\Theta}\mathcal{L}_{\mathrm{rec}}$
\EndWhile
\State $\mathcal{T}
\gets \{\mathcal{C},E_{\phi},
\mathrm{SA},\mathrm{CA},P,\mathrm{FSQ}\}$
\State $\mathcal{T}^{-1}
\gets \{\mathrm{Dequant},\mathrm{Pair},
D_{\theta},\mathcal{C}^{-1}\}$
\State \Return $\mathcal{T},\mathcal{T}^{-1}$
\end{algorithmic}
\end{algorithm}

\paragraph{Bimanual action encoding and quantization.}
We encode each arm's temporal structure with a shared encoder $E_{\phi}$,
a linear projection to width $d$ followed by two causal convolutional
stages:
\begin{equation}
    \mathbf{H}^{r}_{\ell}
    =\operatorname{Mish}\!\left(
    \operatorname{LN}\!\left(
    \operatorname{CConv}_{k_{\ell},s_{\ell}}
    (\mathbf{H}^{r}_{\ell-1})\right)\right),
    \quad \ell=1,2,
    \label{eq:bicat_temporal}
\end{equation}
where $\mathbf{H}^{r}_{0}$ is the projected action sequence. Kernels $(5,3)$
and strides $(2,2)$ compress $H=32$ steps into $M=8$ latent positions.
Each convolution uses left-only padding $k_{\ell}-s_{\ell}$, and layer
normalization acts on features independently at each timestep. Fixed
sinusoidal positions and shared pre-normalized causal self-attention produce
$\mathbf{S}^{L},\mathbf{S}^{R}\in\mathbb{R}^{M\times d}$.
Latent position $i$ thus depends only on actions up to timestep $4i$, using
one-based indexing. The projection, convolutions, normalization, and
self-attention parameters are shared across arms.

We next introduce partner-arm context through symmetric causal
cross-attention. Let $\operatorname{CA}(\mathbf{U},\mathbf{V})$ denote
multi-head attention~\cite{vaswani2017attention} with pre-normalized queries from $\mathbf{U}$ and
pre-normalized keys and values from $\mathbf{V}$. One shared attention block
updates both streams:
\begin{equation}
    \begin{aligned}
        \mathbf{C}^{L}&=\mathbf{S}^{L}
        +\operatorname{CA}(\mathbf{S}^{L},\mathbf{S}^{R}),\\
        \mathbf{C}^{R}&=\mathbf{S}^{R}
        +\operatorname{CA}(\mathbf{S}^{R},\mathbf{S}^{L}).
    \end{aligned}
    \label{eq:bicat_coordination}
\end{equation}
The residual update provides a direct path from each arm's own features while
adding context from its partner. Both directions read the original features
$(\mathbf{S}^{L},\mathbf{S}^{R})$, and query $i$ can attend only to partner
positions $j\leq i$. With dropout disabled, exchanging the
two canonical input streams exchanges their coordinated features. This
symmetry applies to the encoder; the two output streams retain explicit arm
indices and are each conditioned on the partner arm.

A shared linear projection maps each coordinated feature to
$\mathbf{v}^{r}_{i}\in\mathbb{R}^{d_z}$. Finite scalar quantization
(FSQ)~\cite{mentzer2024fsq} then bounds and rounds each coordinate to a fixed
grid, using a straight-through estimator during training. For scalar level
counts $(L_1,\ldots,L_{d_z})$, the grid digits
$n^{r}_{i,j}\in\{0,\ldots,L_j-1\}$ define a categorical symbol:
\begin{equation}
    y^{r}_{i}=\sum_{j=1}^{d_z}n^{r}_{i,j}\prod_{m<j}L_m,
    \qquad K=\prod_{j=1}^{d_z}L_j.
    \label{eq:bicat_quantization}
\end{equation}
We use $d_z=4$ and levels $(8,8,6,5)$, giving $K=1920$. Both arms share the
projection and grid, so the same symbol maps to the same quantized vector in
either stream. We serialize the tokens as
$\mathbf{y}=(y^{L}_{1},\ldots,y^{L}_{M},
y^{R}_{1},\ldots,y^{R}_{M})$, preserving two arm-indexed blocks from a
shared vocabulary.

\paragraph{Joint reconstruction and tokenizer training.}
Joint decoding reconstructs both trajectories from temporally aligned arm
codes. We dequantize each symbol into
$\widehat{\mathbf{z}}^{r}_{i}\in\mathbb{R}^{d_z}$ and concatenate the two
arm vectors at each latent position. A linear projection forms $M$ paired
memory slots:
\begin{equation}
    \mathbf{m}_{i}=\mathbf{W}_{m}
    [\widehat{\mathbf{z}}^{L}_{i}\Vert
    \widehat{\mathbf{z}}^{R}_{i}]+\mathbf{b}_{m}.
    \label{eq:bicat_paired_memory}
\end{equation}
A Transformer decoder $D_{\theta}$ reads this memory from $H$ zero-initialized queries
with fixed sinusoidal positions. A single output head predicts the normalized
bimanual chunk $\widehat{\mathbf{X}}\in\mathbb{R}^{H\times 2d_a}$.
Query self-attention is causal, while cross-attention accesses the complete
paired memory. All $H$ actions are reconstructed in parallel without previous
ground-truth actions as inputs. The decoder therefore reconstructs a complete
chunk from its latents, even though the encoder is causal. Inverse
normalization and the self-inverse right-arm sign transform return
$\widehat{\mathbf{A}}$ to the physical action convention.

We train the tokenizer solely through reconstruction in the shared normalized
space:
\begin{equation}
    \mathcal{L}_{\mathrm{rec}}
    =\mathbb{E}_{\mathbf{A}\sim\mathcal{D}_{\mathrm{valid}}}
    \left[\frac{\|\widehat{\mathbf{X}}-\mathbf{X}\|_{F}^{2}}
    {2Hd_a}\right],
    \label{eq:bicat_tokenizer_loss}
\end{equation}
where $\mathbf{X}=[\mathbf{X}^{L}\Vert\mathbf{X}^{R}]$ and
$\mathcal{D}_{\mathrm{valid}}$ contains chunks with no padded timesteps.
This objective jointly trains the temporal encoder, cross-arm attention,
quantization projection, and decoder, as summarized in
Algorithm~\ref{alg:bicat_training}. FSQ introduces no learnable codebook
or auxiliary commitment loss. We use width $d=256$, eight attention heads,
and dropout $0.1$. The decoder has six pre-normalized layers with feed-forward
width $4d$.

\subsection{Policy Training}
\label{sec:bicat_policy}

After tokenizer training, we freeze \textbf{BiCAT} and use its tokens to
supervise a pretrained VLM backbone with parameters $\theta_{\mathrm{V}}$:
PaliGemma~2-3B~\cite{steiner2024paligemma2} in simulation and
Molmo2-ER~\cite{fang2026molmoact2} on real robots (Sec.~\ref{sec:experiments}). The observation $\mathbf{o}$ forms the input prefix. We augment the vocabulary
with $K$ action symbols and an input-only mask token~\cite{chang2022maskgit}.

\paragraph{Semi-autoregressive action prediction.}
Following semi-autoregressive decoding~\cite{wang2018semiautoregressive},
the policy instantiates Eq.~(\ref{eq:ar_objective}) with the two arm streams
as successive blocks:
\begin{equation}
    p_{\theta_{\mathrm{V}}}(\mathbf{y}\mid\mathbf{o})
    =p_{\theta_{\mathrm{V}}}(\mathbf{y}^{L}\mid\mathbf{o})
    p_{\theta_{\mathrm{V}}}(\mathbf{y}^{R}\mid\mathbf{o},\mathbf{y}^{L}).
    \label{eq:bicat_bar}
\end{equation}
The $M=8$ tokens within each arm block are predicted in parallel, and the
right-arm block is conditioned on the complete left-arm block.
During training, the input suffix
$([\mathrm{MASK}]^{M},\mathbf{y}^{L})$ predicts
$(\mathbf{y}^{L},\mathbf{y}^{R})$ under block-causal attention. This lets the
right block use the complete ground-truth left block without exposing its
own targets. Both blocks receive equally weighted cross-entropy supervision
$\mathcal{L}_{\mathrm{CE}}$ in one forward pass. At inference, two sequential
block predictions generate the action tokens, and the frozen decoder
reconstructs the continuous chunk. We use greedy decoding restricted to
action symbols.
